\documentclass[a4paper,11pt]{article}

\usepackage[margin=1in]{geometry}
\usepackage[T1]{fontenc}
\usepackage[utf8]{inputenc}
\usepackage{lmodern}
\usepackage{microtype}
\usepackage{amsmath,amssymb,mathtools}
\usepackage{booktabs}
\usepackage{hyperref}
\usepackage[parfill]{parskip}

\hypersetup{
  colorlinks=true,
  linkcolor=black,
  citecolor=black,
  filecolor=black,
  urlcolor=black
}

\newcommand{\Zp}{\mathbb Z_p}
\newcommand{\wordeq}{\;\equiv\;}
\newcommand{\rewrites}{\;\Longrightarrow\;}
\newcommand{\xto}{\mathsf X}
\newcommand{\zto}{\mathsf Z}
\newcommand{\sGate}{\mathsf S}
\newcommand{\hGate}{\mathsf H}
\newcommand{\czGate}{\mathsf {CZ}}
\newcommand{\mg}{\mathsf {Mg}}
\newcommand{\mGate}{\mathsf M}
\newcommand{\pauli}{\mathcal P}
\newcommand{\symp}{\mathcal S}
\newcommand{\cliff}{\mathcal C}
\newcommand{\up}{\uparrow}
\newcommand{\down}{\downarrow}
\newcommand{\onehalf}{\tfrac12}
\newcommand{\code}[1]{\texttt{#1}}

\title{Notes on the Agda Development for Qupit Clifford Presentations}
\author{}
\date{\today}

\begin{document}
\maketitle

\section{Overview}

This repository contains the Agda development accompanying the paper
``A Complete and Natural Rule Set for Multi-Qudit Clifford Circuits in
All Odd Prime Dimensions.''  The formalization is organized around
presentations of groups by generators and relations.  Words are built
from gates such as $\sGate$, $\hGate$, and $\czGate$; relations are
encoded as Agda constructors; and derivations are proofs in the
congruence closure generated by those constructors.

There are two intertwined pieces of structure.
First, the symplectic normal-form proof reduces the large family of
box relations to a small set of axioms.  Second, the Clifford
presentation modulo scalars is obtained from a semidirect product of
Pauli translations and symplectic gates.  The semidirect product is
not only a conceptual explanation: it is used in the code to transport
relations from the symplectic and Pauli presentations into the
Clifford presentation.

The most relevant files are:
\[
\begin{array}{ll}
\code{N/Boxes.agda},\ \code{N/Boxes-Sym.agda}
  & \text{box-shaped normal-form indices},\\
\code{N/BR/One},\ \code{N/BR/Two},\ \code{N/BR/Three}
  & \text{box-relation derivations},\\
\code{N/Symplectic-Simplified.agda}
  & \text{small symplectic relation set},\\
\code{N/Clifford/SDProduct.agda}
  & \text{Pauli--symplectic semidirect product},\\
\code{N/Clifford/Iso.agda},\ \code{N/Clifford/Iso3.agda}
  & \text{isomorphism with the Clifford presentation},\\
\code{N/Clifford/Clifford-Mod-Scalar.agda}
  & \text{Clifford modulo scalars}.
\end{array}
\]

\section{How Box Relations Are Reduced}

\subsection{Boxes and normal forms}

The box language packages repeated circuit fragments into finite
families of parameters over $\Zp$.  In \code{N/Boxes.agda} and
\code{N/Boxes-Sym.agda}, the main families are:
\[
\begin{array}{ccl}
A &=& \{(a,b)\in \Zp^2 \mid (a,b)\ne(0,0)\},\\
B &=& \Zp^2,\\
D &=& \Zp^2,\\
E &=& \Zp .
\end{array}
\]
The types $L$, $M$, $LM$, and $NF$ assemble these boxes into
width-indexed normal forms.  Informally, an $NF(n)$ is a triangular
normal form for an $n$-qupit circuit: it consists of an $NF(n-1)$
part followed by the last layer $LM(n)$.  The proof task is then:
when a generator is multiplied against a box normal form, rewrite the
product back into the appropriate box normal form.

The reduction of box relations is deliberately mechanical.  A typical
proof has the following shape.

\begin{enumerate}
\item Expand the named box into the underlying gate word.
\item Use local presentation relations and already-derived pushing
      lemmas to move the incoming generator through the box.
\item Collect the remaining word back into a box constructor.
\item Discharge the parameter update by arithmetic in $\Zp$.
\end{enumerate}

The box-relation files are split by arity and by box family.  For
example, \code{N/BR/One/A.agda} proves one-qupit $A$-box relations,
while \code{N/BR/Two/L-CZ.agda} proves the rules for pushing a
$\czGate$ through a two-qupit $L$ box.  The arithmetic side
conditions used by several proofs are collected in
\code{N/BR/Calculations.agda}.

\subsection{A simple reduction example}

One of the smallest examples is the one-qupit $E$-box relation shown
in the existing literate documentation.  The $E$ box is represented by
a single parameter $b\in\Zp$, and its interpretation is
\[
  [b]^e = \sGate^{-b}.
\]
Multiplying by $\sGate$ on the right gives:
\[
\begin{aligned}
  [b]^e\,\sGate
    &= \sGate^{-b}\sGate \\
    &= \sGate^{-b+1} \\
    &= \sGate^{-(b-1)} \\
    &= [b-1]^e .
\end{aligned}
\]
In Agda this is the proof
\code{lemma-single-qupit-br-E} in \code{doc/Readme.lagda.tex}.  The
only group-theoretic content is the power-combination lemma for
$\sGate$; the rest is modular arithmetic, namely rewriting
$-b+1$ as $-(b-1)$ in $\Zp$.

This example illustrates the pleasant case: the shape of the box does
not change, and the whole relation is a parameter update.

\subsection{A more complicated reduction example}

A more representative example is the general nonzero $A$-box pushed
through $\hGate$.  For $a,b\in\Zp^\times$, the relation proved in
\code{N/BR/One/A.agda} is:
\[
  [(a,b)]^a\,\hGate
  \;\approx\;
  \sGate^{(ab)^{-1}}\;[(b,-a)]^a .
\]
This is the Agda theorem \code{fig-24-3}.  The proof begins by
expanding the $A$ box into the one-qupit normal-form word used by the
development.  After expansion, the core subword is of the form
\[
  \mGate(a^{-1})\,\hGate\,\sGate^{-b/a}\,\hGate .
\]
The derived pushing lemma \code{derived-7} rewrites this into a word
whose parameters are not yet in the desired box format:
\[
  \sGate^{-x^{-1}y^2}\,
  \mGate(-y/x)\,
  \hGate\,\sGate^{-x^{-1}},
  \qquad
  x=-b/a,\quad y=a^{-1}.
\]
The remaining work is arithmetic.  The file
\code{N/BR/Calculations.agda} proves exactly the identities needed to
identify these parameters with the target box:
\[
  -x^{-1}y^2 = (ab)^{-1},\qquad
  -y/x = b^{-1},\qquad
  -x^{-1} = a b^{-1}.
\]
After these substitutions, the word folds back into
\[
  \sGate^{(ab)^{-1}}\;[(b,-a)]^a .
\]

This example shows why the box-relation layer is useful.  The
high-level relation is simple to state, but its proof combines
presentation rewriting, normal-form recognition, nonzero side
conditions, and modular-field calculations.  Agda keeps all of those
bookkeeping steps explicit.

\section{From Symplectic and Pauli Relations to Clifford Relations}

\subsection{The semidirect product presentation}

The Clifford group modulo scalars decomposes as a Pauli translation
part acted on by a symplectic part.  In the code this is made precise
in \code{N/Clifford/SDProduct.agda}.  The generator type is a tagged
sum:
\[
  \mathrm{Gen}_{SD}(n) = \mathrm{Gen}_{XZ}(n) \;\sqcup\;
                         \mathrm{Gen}_{Sym}(n).
\]
The left side contains the Pauli generators $\xto$ and $\zto$; the
right side contains the symplectic gates $\hGate$, $\sGate$, and
$\czGate$.  The relation set is built from three pieces:
\[
  \Gamma_{SD}
  =
  \Gamma_{XZ}
  \;\cup\;
  \Gamma_{Sym}
  \;\cup\;
  \Gamma_{\mathrm{conj}} .
\]
The first two pieces are the Pauli and symplectic presentations.
The third piece records the action of the symplectic generators on
Pauli generators.  For example, \code{conj} in
\code{SDProduct.agda} includes:
\[
\begin{array}{rclcrcl}
\hGate\,\xto\,\hGate^{-1} &=& \zto,
&\qquad&
\hGate\,\zto\,\hGate^{-1} &=& \xto^{-1},\\
\sGate\,\xto\,\sGate^{-1} &=& \xto\zto,
&&
\sGate\,\zto\,\sGate^{-1} &=& \zto,\\
\czGate\,\xto_{\mathrm{top}}\,\czGate^{-1}
  &=& \xto_{\mathrm{top}}\zto_{\mathrm{bottom}},
&&
\czGate\,\xto_{\mathrm{bottom}}\,\czGate^{-1}
  &=& \xto_{\mathrm{bottom}}\zto_{\mathrm{top}}.
\end{array}
\]
These are the familiar Clifford conjugation rules.  In the
semidirect product, they are not derived later as accidental
lemmas; they are the middle relations that say how the right-hand
symplectic factor acts on the left-hand Pauli factor.

\subsection{Embedding the symplectic presentation into Clifford}

The important point is that the symplectic generator $\sGate$ does
not embed as the raw Clifford generator named $\sGate$.  In
\code{N/Clifford/Clifford-Mod-Scalar.agda}, the Clifford presentation
defines derived Pauli words
\[
  \zto = \hGate\hGate\sGate\hGate\hGate\sGate^{-1},
  \qquad
  \xto = \hGate\sGate\hGate\hGate\sGate^{-1}\hGate,
\]
and then defines the symplectic shear
\[
  \mathsf{s} = \sGate\,\zto^{1/2}.
\]
This derived word $\mathsf{s}$ is the image of the symplectic
$\sGate$ in the Clifford presentation.

The forward map from the semidirect product presentation to the
Clifford presentation is defined in \code{N/Clifford/Iso.agda}:
\[
\begin{array}{rcl}
f(\xto) &=& \text{derived Clifford }\xto,\\
f(\zto) &=& \text{derived Clifford }\zto,\\
f(\hGate) &=& \hGate,\\
f(\sGate_{\mathrm{sym}}) &=& \mathsf{s}
      = \sGate_{\mathrm{cliff}}\,\zto^{1/2},\\
f(\czGate) &=& \czGate .
\end{array}
\]
This is the embedding of the symplectic presentation into Clifford:
$\hGate$ and $\czGate$ embed directly, while the symplectic shear
embeds as the Clifford phase gate corrected by a half-$\zto$ Pauli.
That correction is exactly what makes the symplectic action match the
Clifford action modulo scalars in odd prime dimension.

The inverse map explains the correction from the other direction:
\[
\begin{array}{rcl}
h(\hGate_{\mathrm{cliff}}) &=& \hGate,\\
h(\sGate_{\mathrm{cliff}}) &=& \zto^{-1/2}\,\sGate_{\mathrm{sym}},\\
h(\czGate_{\mathrm{cliff}}) &=& \czGate .
\end{array}
\]
The proof \code{f-left-inv-gen} in \code{N/Clifford/Iso3.agda}
checks that these two descriptions cancel, using
$\zto^{-1/2}\zto^{1/2}=1$ and the commutation of $\zto$ with
$\sGate$ in the relevant one-qupit fragment.

\subsection{Deriving Clifford relations}

Once the embedding is fixed, the Clifford relations modulo scalars are
transported from the semidirect product.  The symplectic relations
become Clifford relations after replacing
\[
  \sGate_{\mathrm{sym}}
  \quad\text{by}\quad
  \mathsf{s}=\sGate_{\mathrm{cliff}}\zto^{1/2}.
\]
For example, the symplectic diagonal relation
\[
  \mg\,\sGate = \sGate^{g^2}\,\mg
\]
is transported to the Clifford relation
\[
  \mg\,\mathsf{s} = \mathsf{s}^{g^2}\,\mg,
\]
which appears as \code{semi-M}\ensuremath{\mathsf{s}} in
\code{N/Clifford/Clifford-Mod-Scalar.agda}.  Similarly, the
symplectic $\mGate$ word
\[
  \mGate(x)=\sGate^x\hGate\sGate^{x^{-1}}\hGate\sGate^x\hGate
\]
is sent to the Clifford word with every symplectic $\sGate$ replaced
by $\mathsf{s}$.

The Pauli relations and the mixed action relations give the remaining
Clifford--Pauli transport rules.  For instance, the semidirect action
of $\czGate$ on $\xto$ becomes:
\[
  \czGate\,\xto_{\mathrm{top}}
  =
  \xto_{\mathrm{top}}\zto_{\mathrm{bottom}}\czGate,
  \qquad
  \czGate\,\xto_{\mathrm{bottom}}
  =
  \xto_{\mathrm{bottom}}\zto_{\mathrm{top}}\czGate,
\]
which are the relations \code{rel-X}\ensuremath{\uparrow}\code{-CZ}
and \code{rel-X}\ensuremath{\downarrow}\code{-CZ} in the
Clifford presentation.

The isomorphism theorem in \code{N/Clifford/Iso3.agda} packages this
transport formally:
\[
  \mathcal P \rtimes \mathcal S
  \;\cong\;
  \cliff/\text{scalars}.
\]
Its proof supplies both directions: every semidirect-product relation
is sound in the Clifford presentation, and every Clifford generator is
represented by a semidirect-product word.

\section{Why the Simplified Clifford Relations Have Pauli Tails}

The file \code{N/Clifford/Clifford-Mod-Scalars-Simplified.agda}
tries to state the Clifford relations using the raw gate $\sGate$ as
much as possible, rather than the symplectic shear
$\mathsf{s}=\sGate\zto^{1/2}$.  Most relations are unchanged.  The
interesting cases are the two Selinger $\czGate\hGate\czGate$
relations.

In the purely symplectic presentation, one has for example:
\[
  \czGate\,\hGate_{\mathrm{top}}\,\czGate
  =
  \sGate^{-1}_{\mathrm{top}}\,
  \hGate_{\mathrm{top}}\,
  \sGate^{-1}_{\mathrm{top}}\,
  \czGate\,
  \hGate_{\mathrm{top}}\,
  \sGate^{-1}_{\mathrm{top}}\,
  \sGate^{-1}_{\mathrm{bottom}} .
\]
In the Clifford presentation, however, the embedded symplectic
$\sGate^{-1}$ is $\mathsf{s}^{-1}$, not the raw Clifford
$\sGate^{-1}$.  Replacing $\mathsf{s}^{-1}$ by raw
$\sGate^{-1}$ introduces half-$\zto$ corrections.  These corrections
are pushed through the surrounding $\hGate$ and $\czGate$ gates using
the Pauli action rules.  The corrections then cancel except for a
rightmost Pauli tail.  The simplified relation therefore becomes:
\[
\begin{aligned}
  \czGate\,\hGate_{\mathrm{top}}\,\czGate\,
  \xto_{\mathrm{top}}\zto_{\mathrm{top}}
  &=
  \sGate^{-1}_{\mathrm{top}}\,
  \hGate_{\mathrm{top}}\,
  \sGate^{-1}_{\mathrm{top}}\,
  \czGate\,
  \hGate_{\mathrm{top}}\,
  \sGate^{-1}_{\mathrm{top}}\,
  \sGate^{-1}_{\mathrm{bottom}} .
\end{aligned}
\]
This is \code{selinger-c10} in the simplified Clifford file; the
bottom-qupit analogue is \code{selinger-c11}.  Their soundness is
verified in \code{N/Clifford/Clifford-Simplified-Verify.agda}, and
the identity isomorphism between the original and simplified
Clifford presentations is recorded in
\code{N/Clifford/Iso-Clifford-Simplified.agda}.

\end{document}
